\documentclass[11pt]{article}
\usepackage[a4paper,margin=25mm]{geometry}
\usepackage{amsmath,amssymb,amsthm}
\newtheorem{theorem}{Theorem}\newtheorem{lemma}[theorem]{Lemma}
\title{Literal structural certificates for $M_{22}$ and $M_{22}:2$}
\author{Private L054 finite application}\date{}
\begin{document}\maketitle
\begin{theorem}
Relative to the acquired numerical inventories and named-group interfaces,
all 646 rows of the current companion tables have their displayed structural
descriptions: 156 rows for $M_{22}$ and 490 for $M_{22}:2$.
The descriptions assert existence of the displayed decompositions; they do
not assert uniqueness of a group defined by an extension label.
\end{theorem}
\paragraph{Source and exact models.}
The current companion PDFs are attributed to Thomas Connor and Dimitri
Leemans. All 32 native pages were directly viewed. Powers, parentheses,
colons and dots are preserved, with the published default left association.
For example, row 43 of $M_{22}:2$ is $2^{1+4}\mathbin{\cdot}(2\times6)$,
row 57 is $2^{1+4}\mathbin{\cdot}2^3$, row 179 is
$D_8\mathbin{\cdot}(2\times4)$, and row 180 is
$Q_8\mathbin{\cdot}2^3$. A dot is verified as nonsplitting.
The acquired source bijections and robust numerical inventories identify
all actual subgroup occurrences. Every imported group is given by literal
permutations on the same 22 points; no displayed constructor name is
interpreted as a proper subgroup embedding.

The successful original \texttt{MathieuGroup(22)} execution supplies the
named top model, with its exact driver, clean receipt and all recorded
input/output hashes. Its ordered generators equal the current top
$M_{22}$ generators. The GAP Reference Manual, 50.1-13, is the retained
constructor identification premise. The second named input, $L_3(4)$,
is the accepted L100 projective construction's exact 21-point model:
its retained clean export explicitly checks equality to \texttt{PSL(3,4)}.
All output generators and the complete accepted construction interface
are acquired. Neither name is inferred from its group order.

\paragraph{Registry transport.}
There are 175 earlier accepted structural models and these two named inputs.
The application constructs 97 further models: 58 split extensions,
24 nonsplit extensions, 14 direct products and one quaternion model.
Every occurrence and factor map is a literal bijection fixing identity and
preserving all right-generator edges. The acquired MC.3 word-induction
argument therefore proves that each map is an isomorphism.
For extensions, normality is checked on generators. An actual complement
has trivial intersection with the normal factor and the required product
of orders. Direct factors additionally commute. A displayed nontrivial
semidirect action has a literal noncommuting generator pair.

\newpage
\begin{lemma}[Complete lifts for a finite abelian quotient]
Let $N\mathrel{\trianglelefteq}H$ be finite. Suppose $x_1,\ldots,x_k\in H$
and integers $d_i\geq2$ satisfy $x_i^{d_i}\in N$ and
$[x_i,x_j]\in N$. Suppose moreover that the cosets
\[
 N x_1^{e_1}\cdots x_k^{e_k}\qquad(0\leq e_i<d_i)
\]
are disjoint and cover $H$. Then $H/N\cong\prod_i C_{d_i}$, with the
specified generators. The extension has a complement if and only if there
exist $y_i\in Nx_i$ satisfying $y_i^{d_i}=1$ and $[y_i,y_j]=1$ for all
$i,j$.
\end{lemma}
\begin{proof}
The quotient relations give a surjection from the indicated product of
cyclic groups to $H/N$. The complete distinct normal forms give the same
order, so this map is an isomorphism. If a complement exists, the quotient
map restricted to it is an isomorphism; its inverse supplies the required
lifts. Conversely, lifts satisfying the relations define a homomorphism
from the product of cyclic groups to $H$. Its composition with the quotient
map is the specified isomorphism. It is therefore injective, its image
intersects $N$ trivially, and its image together with $N$ generates $H$.
That image is a complement.
\end{proof}

\paragraph{Exact nonsplitting checks.}
Four new certificates use the lemma, for the actual quotients
$C_2\times C_6$, $C_2^3$ (twice), and $C_2\times C_4$.
The reader first checks every quotient relation, every normal-form coset,
and their disjoint complete union. It then tests every tuple in
$Nx_1\times\cdots\times Nx_k$ against every order and commutation relation.
No satisfying tuple exists. The largest normal subgroup has order 32,
so at most $32^3$ lift tuples occur in any one certificate. Failure of a
candidate search is not a premise. Previously accepted prime, cyclic-six,
$S_3$ and Klein-four criteria pay the other nonsplit descriptions.

The new $Q_{16}$ model uses
$a^8=1$, $b^2=a^4$ and $b^{-1}ab=a^{-1}$, with $a$ of order eight and
$a,b$ generating the actual group of order 16. These relations reduce every
word to $a^i b^j$, $0\leq i<8$, $0\leq j<2$; the actual order proves the
standard quaternion normal forms are distinct. This is the same acquired
dicyclic presentation argument, here at order 16.

\paragraph{Evidence and limits.}
One bounded witness-discovery pass completed all 646 rows in 55.37 seconds,
with clean output. GAP proposed decompositions and maps only. The independent
reader uses literal permutation arithmetic; replay repeats neither discovery
nor any ambient subgroup census. Ordinary and optimized Python must agree.
Small split and nonsplit controls exercise all three new quotient types;
damaged orders, dependent normal forms, missing generators and false dots
are rejected. Source-row, punctuation, degree, embedding and named-model
joins have separate guards, and the earlier complete transport/extension
controls remain acquired.

This application supplies current-table structural joins. It does not claim
historical 2013 PDF byte identity, canonical extension labels, whole-source
admission, new question resolution, formalization or measurement.
Authorship is \texttt{l044\_resume}; root directed registry reuse and
conducts the distinct review.
\end{document}
